\section{Related work}
\paragraph{Text-to-Image Diffusion Models.}
Diffusion models~\cite{ho2020denoising, song2020denoising} have made remarkable progress in text-conditioned image generation~\cite{ saharia2022photorealistic, rombach2022high, kawar2022imagic, ramesh2022hierarchical}, attracting widespread attention in recent years.
The remarkable performance of these models can be attributable to  high-quality large-scale text-image datasets~\cite{changpinyo2021conceptual,schuhmann2021laion,schuhmann2022laion}, the continuous upgrades of foundational models~\cite{chen2023pixart, peebles2023scalable}, conditioning 
encoders~\cite{radford2021learning,ilharco_gabriel_2021_5143773, raffel2020exploring}, and the improvement of controllability~\cite{li2023gligen, zhang2023adding, mou2023t2i, ye2023ip}. 
Due to these advancements,  Podell~\etal~\cite{podell2023sdxl} developed the currently most powerful open-source generative model, SDXL.
Given its impressive capabilities in generating human portraits, we build our PhotoMaker based on this model. 
However, our method can also be extended to other text-to-image synthesis models.




\paragraph{Personalization in Diffusion Models.}
Owing to the powerful generative capabilities  of the diffusion models, more researchers try to explore personalized generation based on them.
Currently, mainstream personalized synthesis methods can be mainly divided into two categories. One relies on additional optimization during the test phase, such as DreamBooth~\cite{ruiz2022dreambooth} and Textual Inversion~\cite{gal2022image}. 
Given that both pioneer works require substantial time for fine-tuning, some studies have attempted to expedite the process of personalized customization by reducing the number of parameters needed for tuning~\cite{kumari2022multi, db_lora, han2023svdiff, yuan2023inserting} or by pre-training with large datasets~\cite{gal2023designing, ruiz2023hyperdreambooth}.
 Despite these advances, they still require extensive fine-tuning of the pre-trained model for each new concept, making the process time-consuming and restricting its applications.
 Recently, some studies~\cite{wei2023elite, ma2023unified, jia2023taming, chen2023subject, shi2023instantbooth, ma2023subject, chen2023anydoor} attempt to perform personalized generation using a single image with a single forward pass, significantly accelerating the personalization process.
These methods either utilize personalization datasets~\cite{chen2023subject, sohn2023styledrop} for training~ or encode the images to be customized  in the semantic space~\cite{wei2023elite, ma2023unified, jia2023taming, 
chen2023photoverse,
shi2023instantbooth, xiao2023fastcomposer}.
Our method focuses on the generation of human portraits based on both of the aforementioned technical approaches.
Specifically,
it not only relies on the construction of an ID-oriented personalization dataset, but also on obtaining the embedding that represents the person's ID in the semantic space.
Unlike previous embedding-based methods, our PhotoMaker extracts a stacked ID embedding from multiple ID images. 
While providing better ID representation, the proposed method can maintain the same high efficiency as previous embedding-based methods.

